% Adiba Shaira - CV
% Compiles on Overleaf with pdfLaTeX (Menu > Compiler > pdfLaTeX). One page, US letter.

\documentclass[10pt, letterpaper]{article}

% Packages:
\usepackage[
    ignoreheadfoot,
    top=1.0 cm,
    bottom=1.0 cm,
    left=1.7 cm,
    right=1.7 cm,
    footskip=1.0 cm,
]{geometry}
\usepackage{titlesec}
\usepackage[dvipsnames]{xcolor}
\definecolor{primaryColor}{RGB}{0, 0, 0}
\usepackage{enumitem}
\usepackage{amsmath}
\usepackage[
    pdftitle={Adiba Shaira - CV},
    pdfauthor={Adiba Shaira},
    pdfcreator={LaTeX},
    colorlinks=true,
    urlcolor=primaryColor
]{hyperref}
\usepackage{bookmark}
\usepackage{needspace}
\usepackage{iftex}

% Ensure machine-readable PDF:
\ifPDFTeX
    \input{glyphtounicode}
    \pdfgentounicode=1
    \usepackage[T1]{fontenc}
    \usepackage[utf8]{inputenc}
\fi

\usepackage{charter}

% Settings:
\raggedright
\pagestyle{empty}
\setcounter{secnumdepth}{0}
\setlength{\parindent}{0pt}
\setlength{\topskip}{0pt}
\pagenumbering{gobble}

\titleformat{\section}{\needspace{4\baselineskip}\bfseries\large}{}{0pt}{}[\vspace{1pt}\titlerule]
\titlespacing{\section}{-1pt}{0.12 cm}{0.08 cm}

\renewcommand\labelitemi{$\vcenter{\hbox{\small$\bullet$}}$}
\newenvironment{highlights}{
    \begin{itemize}[
        topsep=0.06 cm,
        parsep=0.04 cm,
        partopsep=0pt,
        itemsep=0pt,
        leftmargin=10pt
    ]
}{
    \end{itemize}
}

% Entry line: title text flows first, date follows on the same line (correct ATS reading order)
\newcommand{\entry}[2]{\par\noindent #1 \hfill #2\par}

% Publication entry: title, authors, venue
\newcommand{\publication}[3]{\par\noindent\textbf{#1}\par\noindent #2. #3\par}
\newcommand{\me}{\textbf{Adiba Shaira}}

% Separator for the contact lines
\newcommand{\sep}{\kern 5.0 pt$|$\kern 5.0 pt}

\newenvironment{header}{
    \setlength{\topsep}{0pt}\par\kern\topsep\centering\linespread{1.35}
}{
    \par\kern\topsep
}

\let\hrefWithoutArrow\href

\begin{document}

\begin{header}
    \fontsize{25 pt}{25 pt}\selectfont Adiba Shaira

    \vspace{3 pt}

    \normalsize
    \mbox{Stony Brook, NY}\sep%
    \mbox{+1 (934) 451-9631}\sep%
    \mbox{\hrefWithoutArrow{mailto:adibashairarifa@gmail.com}{adibashairarifa@gmail.com}}%
    \\
    \mbox{\hrefWithoutArrow{https://adibashaira.github.io}{adibashaira.github.io}}\sep%
    \mbox{\hrefWithoutArrow{https://linkedin.com/in/adiba-shaira}{linkedin.com/in/adiba-shaira}}\sep%
    \mbox{\hrefWithoutArrow{https://github.com/AdibaShaira}{github.com/AdibaShaira}}%
\end{header}

\vspace{-0.2 cm}

% ---------- Summary ----------
\section{Summary}
Ph.D. candidate in Computer Science specializing in wearable health sensing, multimodal time-series machine learning, and Large Language Model (LLM) applications in digital health. Hands-on experience building end-to-end pipelines from raw physiological signals to interactive health insights.

% ---------- Education ----------
\section{Education}

\entry{\textbf{Ph.D. in Computer Science}, Stony Brook University (Candidate)}{Aug 2023 -- Aug 2027 (expected)}
\entry{\textbf{M.S. in Computer Science}, Stony Brook University, GPA: 3.95/4.00}{Dec 2025}
\begin{highlights}
    \item \href{https://picasso.cs.stonybrook.edu/}{\underline{PiCASSo Lab}}; advised by \href{https://www3.cs.stonybrook.edu/~jain/}{\underline{Prof.\ Shubham Jain}}
    \item \textbf{Research Areas:} Health Sensing, Wearable Computing, Time-Series Modeling, Digital Phenotyping, Machine Learning, Large Language Models (LLMs), Explainable Machine Learning
\end{highlights}
\entry{\textbf{B.Sc. in Computer Science and Engineering}, Bangladesh University of Engineering and Technology}{May 2022}

% ---------- Experience ----------
\section{Experience}

\entry{\textbf{Research Assistant}, Stony Brook University}{May 2024 -- Present}
\begin{highlights}
    \item Built LiftSafe, a machine learning system that predicts unsafe strength-training repetitions one repetition ahead from a single arm-worn Inertial Measurement Unit (IMU), reaching 87\% balanced accuracy in a 24-participant study with EMG ground truth; published at IEEE/ACM CHASE 2026.
    \item Built CareQA, a multi-agent LLM system that answers caregivers' natural-language questions about their autistic child's sleep, activity, and physiological signals using passively sensed data from the \href{https://www.sfari.org/resource/simons-sleep-project/}{\underline{Simons Sleep Project}}.
    \item Now extending CareQA with an interactive visualization dashboard that pairs a child's sensor trends with LLM-generated explanations, and a caregiver study that grounds the system in caregivers' own questions.
    \item Designing passive smartphone-based sleep sensing for PTSD monitoring in a World Trade Center 9/11 first-responder cohort: extracting sleep metrics from accelerometer and screen-state data and correlating them with clinical PCL scores via longitudinal statistical modeling.
\end{highlights}

\entry{\textbf{Teaching Assistant}, Stony Brook University: \href{https://www3.cs.stonybrook.edu/~skiena/373/}{\underline{CSE 373, Analysis of Algorithms}}}{Aug 2023 -- May 2024}

\entry{\textbf{Lecturer}, United International University, Dhaka, Bangladesh}{May 2022 -- July 2023}
\begin{highlights}
    \item Taught Human-Computer Interaction (HCI), Analysis of Algorithms, and Computer Networks.
\end{highlights}

% ---------- Publications ----------
\section{Publications}

\publication{CareQA: Multi-Agent Sensemaking to Answer Caregiver Queries from Passive Sensing Data in Autism Care}
{\me, Michelle S. Ballan, Shubham Jain}
{\textit{Under submission at IEEE PerCom 2027.}}

\vspace{0.08 cm}
\publication{\href{https://picasso.cs.stonybrook.edu/bibcite/reference/47}{LiftSafe: Predicting Unsafe Repetitions in Strength Training using a Single Wearable Sensor}}
{\me, Hunter J. Bennett, Shubham Jain}
{\textit{IEEE/ACM Conference on Connected Health: Applications, Systems, and Engineering Technologies (CHASE)}, 2026.}

\vspace{0.08 cm}
\publication{\href{https://dl.acm.org/doi/10.1007/978-3-031-58072-7_6}{Gene Tree Parsimony in the Presence of Gene Duplication, Loss, and Incomplete Lineage Sorting}}
{Prottoy Saha, Md.\ Shamiul Islam, Tasnim Rahman, \me, Kazi Noshin, Rezwana Reaz, Md.\ Shamsuzzoha Bayzid}
{\textit{RECOMB International Workshop on Comparative Genomics}, Springer Nature, pp.\ 110--128, 2024.}

% ---------- Other Projects ----------
\section{Other Projects}

\entry{\textbf{\href{https://github.com/jamaliasultana/sleep_anomaly}{Using a Single Device to Predict Sleep Efficiency and Next-Day's Mood}}}{2024}
Modeled next-day sleep efficiency and mood from single-wearable time-series data in the GLOBEM dataset. \textit{Key Techniques:} Python, Pandas, NumPy, Scikit-learn, LSTM, Autoencoder models (PyTorch/TensorFlow)

\vspace{0.08 cm}
\entry{\textbf{\href{https://github.com/shahrarfatemi/global_disaster_visualizer}{Global Disaster Visualizer}}}{2023}
Interactive web-based visualization of global disaster data. \textit{Key Techniques:} D3.js, JavaScript, React, Python

% ---------- Awards ----------
\section{Awards}

\entry{\textbf{NSF Student Travel Award}, IEEE/ACM CHASE}{2026}
\entry{\textbf{Distinguished Travel Award}, Graduate Student Organization, Stony Brook University}{2026}

% ---------- Technical Skills ----------
\section{Technical Skills}

\textbf{Languages:} Python, Java, C/C++, SQL, JavaScript

\textbf{Machine Learning \& Deep Learning:} PyTorch, TensorFlow, Keras, Scikit-learn, Hugging Face Transformers, SHAP, Large Language Models (LLMs), Multi-Agent LLM Systems, Retrieval-Augmented Generation (RAG)

\textbf{Sensing \& Data:} IMU data processing, physiological signals (PPG, EDA, EEG), OpenCV, MATLAB, Pandas, NumPy

\textbf{Tools:} Git, Jupyter, Wireshark, Django, Bootstrap, MySQL, MongoDB

\end{document}
